\chapter{Conception}\label{chap:conception}

\section{Acteurs et besoins}
Trois profils utilisent l'application. Le tableau~\ref{tab:acteurs} résume ce que chacun peut faire.

\begin{table}[htbp]
\centering
\caption{Acteurs et fonctionnalités}\label{tab:acteurs}
\begin{tabularx}{\linewidth}{@{}lX@{}}
\toprule
\textbf{Acteur} & \textbf{Fonctionnalités}\\
\midrule
Visiteur & Parcourir la vitrine, filtrer par catégorie, rechercher, ouvrir la fiche d'un produit, créer un compte, se connecter.\\
Client & Tout ce que fait le visiteur, plus : ajouter un produit au panier, modifier une quantité, retirer une ligne, vider le panier.\\
Administrateur & Créer, modifier, retirer de la vente et supprimer des produits ; gérer les catégories.\\
\bottomrule
\end{tabularx}
\end{table}

\section{Diagramme de classes}
La figure~\ref{fig:classes} (page suivante) donne le modèle du domaine. Les relations suivent la logique du sujet.

\begin{itemize}
  \item \textbf{Héritage} : \code{Client} et \code{Administrateur} spécialisent la classe abstraite \code{Internaute}.
  \item \textbf{Client -- Panier}, 1 à 1 : un client possède un seul panier, créé en même temps que lui. Un
        administrateur n'en a pas.
  \item \textbf{Panier -- LignePanier}, composition : une ligne n'existe que dans son panier ; supprimer le panier
        supprime ses lignes.
  \item \textbf{LignePanier -- Produit}, association : la ligne référence un produit mais ne le possède pas.
  \item \textbf{Vitrine -- Catégorie -- Produit}, agrégations : la vitrine présente des catégories, qui regroupent
        des produits.
  \item \textbf{Administrateur -- Produit}, dépendance : l'administrateur gère le catalogue, sans lien persistant.
\end{itemize}

\begin{landscape}
\thispagestyle{empty}
\centering
\resizebox{\linewidth}{!}{\begin{tikzpicture}[
  % UML class box: title part (tinted) + attribute/method part
  cls/.style={
    rectangle split, rectangle split parts=2,
    draw=black!70, thick, fill=white,
    rectangle split part fill={teal!15, white},
    font=\footnotesize, text width=5.3cm, align=left,
    inner xsep=6pt, inner ysep=4pt},
  abs/.style={cls, draw=teal!60!black},
  mult/.style={font=\scriptsize, fill=white, inner sep=1pt},
  rel/.style={font=\scriptsize\itshape, fill=white, inner sep=1.5pt}
]

% --- Internaute hierarchy ---
\node[abs] (int) at (-2.9,0) {
  \textbf{\textit{Internaute}} \\ {\scriptsize\textit{<<abstract>>}}
  \nodepart{second}
  - id : Long\\
  - nom, prenom : String\\
  - email : String (unique)\\
  - motDePasseHash : String\\
  - dateInscription : LocalDateTime};

\node[cls] (adm) at (-6.0,-5.4) {
  \textbf{Administrateur}
  \nodepart{second}
  - matricule : String};

\node[cls] (cli) at (0.2,-5.4) {
  \textbf{Client}
  \nodepart{second}
  - adresse : String\\
  - telephone : String};

% --- Cart ---
\node[cls] (pan) at (8.8,-5.4) {
  \textbf{Panier}
  \nodepart{second}
  - id : Long\\
  - dateCreation : LocalDateTime\\
  - dateModification : LocalDateTime\\
  + ajouterProduit(Produit, int)\\
  + modifierQuantite(Produit, int)\\
  + retirerProduit(Produit)\\
  + vider()\\
  + getTotal() : BigDecimal};

\node[cls] (lig) at (17.4,-5.4) {
  \textbf{LignePanier}
  \nodepart{second}
  - id : Long\\
  - quantite : int\\
  - prixUnitaire : BigDecimal\\
  + getSousTotal() : BigDecimal};

% --- Catalogue ---
\node[cls] (pro) at (17.4,-11.6) {
  \textbf{Produit}
  \nodepart{second}
  - id : Long\\
  - reference : String (unique)\\
  - nom, description : String\\
  - prix : BigDecimal\\
  - stock : int\\
  - imageUrl : String\\
  - actif : boolean\\
  - version : long (@Version)};

\node[cls] (cat) at (8.8,-11.6) {
  \textbf{Categorie}
  \nodepart{second}
  - id : Long\\
  - nom : String\\
  - description : String};

\node[cls] (vit) at (0.2,-11.6) {
  \textbf{Vitrine}
  \nodepart{second}
  - id : Long\\
  - titre : String\\
  - description : String\\
  - active : boolean};

% --- Inheritance (hollow triangle) ---
\draw[-{Triangle[open,length=8pt,width=8pt]}, thick] (adm.north) -- (int.south);
\draw[-{Triangle[open,length=8pt,width=8pt]}, thick] (cli.north) -- (int.south);

% --- Client 1 - 1 Panier ---
\draw[thick] (cli.east) -- (pan.west)
  node[mult, pos=0.12, above] {1}
  node[mult, pos=0.88, above] {1}
  node[rel, midway, below] {poss\`ede};

% --- Panier composition LignePanier ---
\draw[{Diamond[length=9pt,width=6pt,fill=black]}-, thick] (pan.east) -- (lig.west)
  node[mult, pos=0.12, above] {1}
  node[mult, pos=0.88, above] {0..*}
  node[rel, midway, below] {contient};

% --- LignePanier -> Produit ---
\draw[-{Stealth[length=7pt]}, thick] (lig.south) -- (pro.north)
  node[mult, pos=0.12, right] {0..*}
  node[mult, pos=0.88, right] {1}
  node[rel, midway, right] {r\'ef\'erence};

% --- Categorie aggregation Produit ---
\draw[{Diamond[length=9pt,width=6pt,fill=white]}-, thick] (cat.east) -- (pro.west)
  node[mult, pos=0.12, above] {1}
  node[mult, pos=0.88, above] {0..*}
  node[rel, midway, below] {regroupe};

% --- Vitrine aggregation Categorie ---
\draw[{Diamond[length=9pt,width=6pt,fill=white]}-, thick] (vit.east) -- (cat.west)
  node[mult, pos=0.12, above] {1}
  node[mult, pos=0.88, above] {0..*}
  node[rel, midway, below] {pr\'esente};

% --- Administrateur manages Produit (dependency, dashed) ---
\coordinate (dep1) at (-6.0,-15.2);
\coordinate (dep2) at (17.4,-15.2);
\draw[-{Stealth[length=7pt]}, thick, dashed]
  (adm.south) -- (dep1) -- (dep2) -- (pro.south);
\node[rel] at (5.7,-15.2) {g\`ere};

\end{tikzpicture}
}
\captionof{figure}{Diagramme de classes de la version A (panier, vitrine, internaute)}\label{fig:classes}
\end{landscape}

\section{Règles métier}
Les règles qui protègent la cohérence du panier sont écrites \textbf{dans l'entité \code{Panier}} et non dans les
servlets ou les services. Une règle placée dans le domaine ne peut pas être oubliée par un nouvel écran : quel que
soit l'appelant, la même vérification s'applique. Le tableau~\ref{tab:regles} les liste avec l'endroit où elles sont
imposées une seconde fois, en base, quand c'est possible.

\begin{table}[htbp]
\centering
\caption{Règles métier et double protection}\label{tab:regles}
\begin{tabularx}{\linewidth}{@{}lXl@{}}
\toprule
\textbf{N°} & \textbf{Règle} & \textbf{Garde-fou en base}\\
\midrule
R1 & Une quantité de ligne est comprise entre 1 et 99. & \code{CHECK quantite >= 1}\\
R2 & Seul un produit actif peut être ajouté au panier. & ---\\
R3 & La quantité d'une ligne ne dépasse pas le stock disponible. & \code{CHECK stock >= 0}\\
R4 & Ajouter un produit déjà présent fusionne la ligne (les quantités s'additionnent). & \code{UNIQUE (panier, produit)}\\
R5 & Le prix de la ligne est figé au moment de l'ajout. & \code{CHECK prix\_unitaire >= 0}\\
R6 & Le total est la somme des sous-totaux, calculée en \code{BigDecimal}. & ---\\
R7 & L'adresse électronique est unique, normalisée en minuscules. & \code{UNIQUE (email)}\\
R8 & La référence d'un produit est unique, en majuscules. & \code{UNIQUE (reference)}\\
R9 & Un produit présent dans un panier ne peut pas être supprimé (on le retire de la vente). & clé étrangère\\
R10 & Une catégorie qui contient des produits ne peut pas être supprimée. & clé étrangère\\
R11 & Un internaute a obligatoirement un rôle (\code{CLIENT} ou \code{ADMIN}). & \code{NOT NULL} + \code{CHECK}\\
\bottomrule
\end{tabularx}
\end{table}

\section{Choix de la stratégie d'héritage}
JPA propose trois façons de stocker une hiérarchie. Le tableau~\ref{tab:heritage} les compare pour notre cas.

\begin{table}[htbp]
\centering
\caption{Stratégies d'héritage JPA}\label{tab:heritage}
\begin{tabularx}{\linewidth}{@{}lXX@{}}
\toprule
\textbf{Stratégie} & \textbf{Avantages} & \textbf{Inconvénients}\\
\midrule
\code{SINGLE\_TABLE} & Une seule table, aucune jointure, authentification par une requête simple. & Colonnes propres aux sous-classes nullables ; une colonne discriminante.\\
\code{JOINED} & Modèle normalisé, pas de colonnes inutiles. & Jointure à chaque lecture, plus de tables.\\
\code{TABLE\_PER\_CLASS} & Chaque classe a sa table complète. & Requêtes polymorphes par \code{UNION}, clés à gérer entre tables.\\
\bottomrule
\end{tabularx}
\end{table}

J'ai retenu \textbf{\code{SINGLE\_TABLE}}. La hiérarchie est petite (deux sous-classes qui ajoutent chacune un
ou deux champs), et la lecture la plus fréquente est la connexion : retrouver un internaute par son adresse sans
savoir à l'avance s'il est client ou administrateur. Avec une table unique, c'est une seule requête sans jointure.
Le défaut de la stratégie, la colonne discriminante laissée nullable par EclipseLink, est corrigé par une contrainte
\code{NOT NULL} et un \code{CHECK} (chapitre~\ref{chap:persistance}).
